% Seitenränder
\geometry{
	paper=a4paper,
	right=20mm,		
	top=20mm,
	left=30mm,
	bottom=30mm
}

% Zeilenabstand 1,5 Zeilen
\singlespacing
%\spacing{1.21}
%\onehalfspacing
%\doublespacing


% Minimal mehr platz nach einem punkt. <-
\frenchspacing

% Schusterjungen und Hurenkinder vermeiden
\clubpenalty = 10000
\widowpenalty = 10000 
\displaywidowpenalty = 10000


% Style der listings
\definecolor{RoyalBlue}{cmyk}{1, 0.50, 0, 0}
\lstset{
	basicstyle=\ttfamily,
	numberstyle=\tiny\color{gray},
	keywordstyle=\color{RoyalBlue},
	commentstyle=\color{green}\ttfamily,
	stringstyle=\color{mauve},
	autogobble=true,
	breaklines=true,
	breakatwhitespace=true,
	captionpos = b,
	frame=single, % box um das listing
	numbers=left, %nummerierung um das Listing
	% Nummerierung innerhalb der box
	framexleftmargin=5.5mm,
	xleftmargin=19pt,
	xrightmargin=3.4pt,
} 


%Ligaturen verhindern:
%\DisableLigatures{}
%Was sind ligaturen ?:
%ff und fi ragen ineinander, und auch andere wilde Buchstabenkombinationen


% Kopf und Fußzeile
\pagestyle{scrheadings}
% Kopf- und Fußzeile auch auf Kapitelanfangsseiten
\renewcommand*{\chapterpagestyle}{scrheadings}
% Schriftform der Kopfzeile
\renewcommand{\headfont}{\normalfont}

% Höhe der Kopfzeile
\setlength{\headheight}{14mm}
% Kopfzeile
\ihead{\headmark}
\chead{}
%\ohead{\includegraphics[scale=0.15, trim={2.4cm 24.5cm 2.6cm 2.6cm}, clip]{Bilder/SIEMENS-Logo.pdf}}
\ohead{\includegraphics[scale=0.08]{Bilder/Siemens_Logo.pdf}}
%\ohead{\art} % Masterarbeit statt SIEMENS in der Kopfzeile

% Fußzeile
% Leer, da sich das Layout ändert
\cfoot{}
\ifoot{}
\ofoot{}

% Layout Inhaltsverzeichnis
\setcounter{tocdepth}{2}  % Inhaltsverzeichnis nur bis 2 Unterelemente auflisten
\setcounter{secnumdepth}{3} % Numerierungstiefe für Inhaltsverzeichnis

\DeclareTOCStyleEntries[
% numwidth=1em,					% Abstand Zahl - Überschrift
rightindent=4em,				% Rand nach Links. Relevant für Überschriften mit Zeilenumbruch im Inhaltsverzeichnis
pagenumberbox=\pagenumberbox
]
{tocline}{chapter, section, subsection, subsubsection, paragraph, subparagraph, figure, table}
\newcommand*\pagenumberbox[1]{\mbox{\hspace{2pt}#1}} 

%Höhe von Kapiteln im Text					2.3
%\renewcommand*{\chapterheadstartvskip}{\vspace*{1.0\baselineskip}}% Abstand einstellen
\renewcommand*{\chapterheadstartvskip}{\vspace*{1\baselineskip}}% Abstand einstellen

% Eigenes Design für Abkürzungsverzeichnis
\newglossarystyle{mylong}{%
	\renewenvironment{theglossary}%
	{\begin{longtable}[l]{llp{\glsdescwidth}p{\glspagelistwidth}}}%
		{\end{longtable}}%
	\renewcommand*{\glossaryheader}{}%
	\renewcommand*{\glsgroupheading}[1]{}%
	\renewcommand{\glossentry}[2]{%
		\textsf{\textbf{\glsentryitem{##1}\glstarget{##1}{\glossentryname{##1}}}} &		\glossentrysymbol{##1} &		\glossentrydesc{##1} &		##2\tabularnewline
	}%
	\renewcommand{\subglossentry}[3]{%
		&
		\glssubentryitem{##2}%
		\glossentrysymbol{##2} &
		\glstarget{##2}{\strut}\glossentrydesc{##2} & ##3\tabularnewline
	}%
	\renewcommand*{\glsgroupskip}{%
		\ifglsnogroupskip\else & & &\tabularnewline\fi}%
}

% Abstand zwischen Bild und Bildunterschrift
% \setlength\abovecaptionskip{-20pt}
%Vektorbasierte Standardschrift
\renewcommand{\familydefault}{\sfdefault} 	% %Defaultschrift ohne Serifen (modernen)
\usepackage{sansmath}						% %Matheschrift ohne Serifen
\usepackage{microtype}						% %Verbesserte Textverteilung in Zeile

